\documentclass[11pt, a4paper]{article}
\usepackage[parfill]{parskip}
\usepackage[utf8]{inputenc}
\usepackage[margin=1in]{geometry}
\usepackage{amsmath}
\usepackage{amssymb}
\usepackage{graphicx}
\usepackage{tabularx}
\usepackage{hyperref}
\usepackage{charter}

\title{\textbf{IM0802 \textendash{} Responsible Artificial Intelligence} \\ Proposed established resources for assignment 1}
\author{Harman Preet Singh (852863316) \\ harman.pnahal@gmail.com}


\begin{document}

\maketitle

\section{Chosen topic: {\normalfont Deepfakes in time of conflict or war}}

% file:///C:/Users/Harman/Desktop/UNI/MS%20Artificial%20Intelligence/JAAR%201/SEM%202/Responsible%20Artificial%20Intelligence/Assignment%201/RAI_Assignment_1_Description_2025-2026.pdf

\section{Introduction: {\normalfont Fictional setting that illustrates the context}}

\textbf{Scenario}: The Fabricated War Crime
A deepfake video depicting soldiers from one nation massacring civilians in a disputed region goes viral hours before a UN Security Council vote. The footage shifts international opinion and triggers sanctions — all based on footage that never happened.

\subsection{Paragraph 1: {\normalfont The scenario}}
Present the scenario as your fictional but plausible ethical dilemma. Write it in a slightly narrative tone to draw the reader in — this is meant to capture attention. Describe the setting (a disputed conflict zone, a live UN vote), the action (the deepfake video going viral), and the impact (sanctions imposed on false grounds). Keep it concise, around 150–200 words.

\subsection{Paragraph 2: {\normalfont The aim of the essay}}
Transition from the scenario into what the essay will actually do. Briefly state the topic (deepfakes in conflict), why it matters ethically, and what the essay will cover — the issue, the stakeholders affected, and the solutions considered. This grounds the reader after the narrative opening and signals the academic intent. Around 100–150 words is plenty here.


\end{document}